\subsection{Orthogonal Decomposition and Its Consequences}

The projection formula splits $\mathbf u$ into two geometrically different parts:
\[
\mathbf u
=
\operatorname{proj}_{\mathbf v}\mathbf u
+
\left(\mathbf u-\operatorname{proj}_{\mathbf v}\mathbf u\right).
\]
The first term is parallel to $\mathbf v$ and the second is perpendicular to
$\mathbf v$. This is the orthogonal decomposition of $\mathbf u$ relative to
the direction of $\mathbf v$.

Because the two components are perpendicular, the Pythagorean theorem gives
\[
\boxed{
\|\mathbf u\|^2
=
\left\|\operatorname{proj}_{\mathbf v}\mathbf u\right\|^2
+
\left\|\mathbf u-\operatorname{proj}_{\mathbf v}\mathbf u\right\|^2
}.
\]

\subsubsection*{Scalar component and vector projection}

The signed scalar component of $\mathbf u$ in the direction of nonzero
$\mathbf v$ is
\[
\boxed{
\operatorname{comp}_{\mathbf v}\mathbf u
=
\frac{\mathbf u\cdot\mathbf v}{\|\mathbf v\|}
}.
\]
It measures how much signed length of $\mathbf u$ lies along the direction of
$\mathbf v$. The vector projection is obtained by multiplying this number by
the unit vector in that direction:
\[
\operatorname{proj}_{\mathbf v}\mathbf u
=
\operatorname{comp}_{\mathbf v}\mathbf u
\frac{\mathbf v}{\|\mathbf v\|}.
\]

\subsubsection*{The Cauchy--Schwarz inequality}

The length of the parallel component cannot exceed the length of the entire
vector:
\[
\left|\operatorname{comp}_{\mathbf v}\mathbf u\right|
\leq \|\mathbf u\|.
\]
Substitution gives
\[
\boxed{
|\mathbf u\cdot\mathbf v|
\leq
\|\mathbf u\|\,\|\mathbf v\|
}.
\]
Therefore
\[
-1\leq
\frac{\mathbf u\cdot\mathbf v}
{\|\mathbf u\|\,\|\mathbf v\|}
\leq1,
\]
which guarantees that the angle formula using $\arccos$ is meaningful.

\subsubsection*{The triangle inequality}

The direct displacement $\mathbf u+\mathbf v$ cannot be longer than travelling
first by $\mathbf u$ and then by $\mathbf v$. Algebraically,
\[
\begin{aligned}
\|\mathbf u+\mathbf v\|^2
&=\|\mathbf u\|^2+2\mathbf u\cdot\mathbf v+\|\mathbf v\|^2\\
&\leq\|\mathbf u\|^2+2\|\mathbf u\|\,\|\mathbf v\|+\|\mathbf v\|^2\\
&=(\|\mathbf u\|+\|\mathbf v\|)^2.
\end{aligned}
\]
Hence
\[
\boxed{
\|\mathbf u+\mathbf v\|
\leq
\|\mathbf u\|+\|\mathbf v\|
}.
\]
This is the algebraic form of the geometric fact that the straight segment is
the shortest route between two points.